\section{Evaluation Protocol}
\label{app:evaluation-protocol}

\textbf{Bits-per-byte implementation.} We measure the negative log-likelihood of held-out text in bits, normalized by its UTF-8 byte count.
For language $\ell$, let $x_1,\ldots,x_{m_\ell}$ be the evaluated token sequence, $c_t$ the context preceding token $x_t$ within the scoring block, and $b_\ell$ the number of bytes in the evaluated text. We compute
\begin{equation}
\mathrm{BPB}_\ell = \frac{-\sum_{t=2}^{m_\ell}\log_2 p(x_t \mid c_t)}{b_\ell}.
\label{eq:bpb}
\end{equation}

We concatenate validation documents without end-of-document separators and score them in blocks with 4,096 input tokens and one-token overlap between successive blocks. Each token except the first of the stream is scored once; context resets at block boundaries. We select a deterministic prefix ending at the document boundary that reaches 1M tokens, or the full stream if it is shorter. We measure the denominator by decoding the entire selected prefix and counting its UTF-8 bytes, including the initial conditioning token. Thus, the implementation omits the first token's likelihood but retains its bytes. We also compute token-level perplexity from the same forward passes.

% The original draft reports a standard error of about 0.003 bits for this
% prefix. Add the estimation procedure before presenting it as uncertainty.

% BEGIN generated: benchmarks (make_appendix_tables.py)
\begin{table*}[t]
\centering
\small
\setlength{\tabcolsep}{4pt}
\resizebox{\textwidth}{!}{
\begin{tabular}{llllrrr}
\toprule
\textbf{Benchmark} & \textbf{Task} & \textbf{Format} & \textbf{Construction} & \textbf{Langs} & \textbf{Trained} & \textbf{Tasks} \\
\midrule
\multicolumn{7}{l}{\textit{multilingual}} \\
Belebele & reading comprehension & 4-way MC & HT (FLORES) & 94 & 49 & 59 \\
Global PIQA & physical commonsense & 2-way completion & native & 77 & 48 & 65 \\
MultiBLiMP & grammatical acceptability & minimal pairs & auto (UD, UniMorph) & 57 & 34 & 34 \\
INCLUDE & regional knowledge & 4-way MC & native exams & 43 & 36 & 36 \\
INCLUDE v2 & regional knowledge & 4-way completion & native exams & 42 & 42 & 77 \\
Global-MMLU & world knowledge & 4-way MC & HT + community & 37 & 29 & 29 \\
HellaSwag (Okapi) & activity commonsense & 4-way completion & MT & 31 & 26 & 26 \\
ARC (Okapi) & science QA & 4-way MC & MT & 32 & 27 & 28 \\
ARC-Challenge MT & science QA & 4-way MC & MT & 11 & 11 & 11 \\
XNLI & natural language inference & 3-way & HT & 18 & 15 & 15 \\
XStoryCloze & narrative commonsense & 2-way ending & HT & 13 & 8 & 8 \\
XCOPA & causal commonsense & 2-way & HT & 11 & 8 & 8 \\
PAWS-X & paraphrase identification & 2-way & HT & 10 & 8 & 8 \\
XWinograd & coreference & Winograd schemas & native & 6 & 6 & 6 \\
LAMBADA multilingual & last-word prediction & completion & MT & 5 & 5 & 5 \\
TruthfulQA-Multi (mc1) & truthfulness & MC & HT & 3 & 2 & 2 \\
TruthfulQA (mc2) & truthfulness & MC, probability mass & native + MT & 3 & 3 & 3 \\
BLEnD (sample) & everyday cultural knowledge & MC & native & 4 & 4 & 5 \\
\midrule
\multicolumn{7}{l}{\textit{English text}} \\
INCLUDE v2 (English) & regional knowledge & 4-way completion & native, translated to en & 42 & 42 & 77 \\
CulturalBench (easy) & cultural knowledge & 4-way MC & native (English) & 8 & 8 & 19 \\
CulturalBench (hard) & cultural knowledge & 2-way (true/false) & native (English) & 8 & 8 & 19 \\
MMLU & world knowledge & 4-way MC & native & 1 & 1 & 1 \\
CommonsenseQA & commonsense & 5-way MC & native & 1 & 1 & 1 \\
OpenBookQA & science QA & 4-way MC & native & 1 & 1 & 1 \\
MathQA & math word problems & 5-way MC & native & 1 & 1 & 1 \\
ToxiGen & toxicity detection & 2-way & auto (LLM-generated) & 1 & 1 & 1 \\
BBH (letter subtasks) & reasoning & letter MC & native & 1 & 1 & 17 \\
BBH (two-way subtasks) & reasoning & 2-way completion & native & 1 & 1 & 6 \\
ACP-Bench (MC) & planning & 4-way letter MC & auto (PDDL) & 1 & 1 & 7 \\
ACP-Bench (boolean) & planning & yes/no completion & auto (PDDL) & 1 & 1 & 7 \\
\midrule
\multicolumn{7}{l}{\textit{reformulated variants}} \\
Belebele (RF) & reading comprehension & cloze completion & HT (FLORES) & 94 & 49 & 59 \\
Belebele (LLM-RF) & reading comprehension & statement continuation & HT (FLORES) & 49 & 49 & 59 \\
INCLUDE (RF) & regional knowledge & cloze completion & native exams & 43 & 36 & 36 \\
INCLUDE (LLM-RF) & regional knowledge & statement continuation & native exams & 43 & 36 & 36 \\
Global-MMLU (RF) & world knowledge & cloze completion & HT + community & 37 & 29 & 29 \\
CulturalBench (easy) (RF) & cultural knowledge & cloze completion & native (English) & 8 & 8 & 19 \\
MMLU (RF) & world knowledge & cloze completion & native & 1 & 1 & 1 \\
CommonsenseQA (RF) & commonsense & cloze completion & native & 1 & 1 & 1 \\
BBH (letter subtasks) (RF) & reasoning & cloze completion & native & 1 & 1 & 17 \\
ACP-Bench (MC) (RF) & planning & cloze completion & auto (PDDL) & 1 & 1 & 7 \\
\midrule
Total & & & & & 50 & 846 \\
\bottomrule
\end{tabular}
}
\caption{Benchmarks evaluated during pretraining. Langs is the number of languages the harness registers for the benchmark; Trained how many of them the sweep trains (English and the FineWeb-2 languages of every trained build); Tasks the number of language tasks evaluated on those languages. A model is evaluated only on the tasks of the languages in its training mixture. The reformulated variants pose the same items as their original: RF scores the answer strings instead of the option letters, LLM-RF rewrites each item into a statement stem with short continuations. MC = multiple choice; QA = question answering; UD = Universal Dependencies. Construction: HT = human or professional translation, MT = machine translation, native = written in the language, auto = automatically generated.}
\label{tab:benchmark-families}
\end{table*}
% END generated: benchmarks

\textbf{Task selection.} We select the registered tasks whose language is in the training mixture; task counts can exceed family counts. For the resource-ranked lists, the number of tasks grows with $L$: 107 at $L=1$, 127 at $L=2$, 306 at $L=8$, 456 at $L=15$, 639 at $L=30$, and 846 at $L=50$. These counts cover the whole evaluated group of Table~\ref{tab:benchmark-families}: the multilingual families, the reformulated variants of the letter-format families (\S\ref{subsec:rq1}) and the English-text benchmarks. Every trained language is covered by at least one family. For the seed-1904 deep scheme-A ladder we additionally run every task in every language, independent of the training mixture, to read cross-lingual transfer on the benchmarks as well as on bits-per-byte. 

\textbf{Language coverage.} Every trained language is covered by at least one benchmark family, and English by 14 of the multilingual families plus 11 of the English-text benchmarks. Of the 100 languages of Table~\ref{tab:app-languages}, eighteen have exactly one family and fourteen have two; the languages with a single family are covered by Belebele (12), MultiBLiMP (4) or Global PIQA (2) alone, and for those the per-language benchmark read rests on one task. Counting tasks rather than families, a cell is evaluated on 107 tasks at $L=1$, 127 at $L=2$, 306 at $L=8$, 456 at $L=15$, 639 at $L=30$ and 846 at $L=50$, the reformulated variants and the English-only benchmarks included. Two IrokoBench families (AfriXNLI and AfriMMLU) are wired but disabled because of a loading defect in the pinned harness; re-enabling them would add a second or third human-translated family for Amharic, Igbo, Kinyarwanda, Southern Sotho, Xhosa and Zulu.


\textbf{Harness settings.} All benchmarks run through the Swiss AI fork of the LM Evaluation Harness \citep{gao_framework_2021}, installed from a pinned shared checkout with the source commit recorded in each evaluation job, with the vLLM backend \citep{kwon_efficient_2023}, a maximum context of 4,096 tokens, a beginning-of-sequence token prepended to every prompt, no chat template, and each task's default prompt and number of in-context examples. We report accuracy for multiple-choice and completion tasks and exact match for the generative ones, and aggregate sub-tasks (for example the MMLU subjects) to their family by the harness's own grouping. Base models at these sizes are at chance on several families. The above-chance threshold of Appendix~\ref{app:chance-threshold} decides per task and size whether such a cell enters the analysis.

